\documentclass[border=4pt]{standalone}
\usepackage{amsmath,amssymb,mathtools,xcolor,varwidth}
\begin{document}
\begin{varwidth}{28em}
\large\bfseries
Here is a single ellipse with centre $O$, drawn wider than tall. Through the centre run two conjugate diameters $PQ$ and $RS$ --- note that they are \emph{skewed}, not perpendicular, for that is the whole point of the lemma. About the ellipse is described the parallelogram whose four sides are tangent to the curve exactly at the four diameter endpoints $P$, $Q$, $R$, $S$. Apollonius proved that no matter which pair of conjugate diameters one chooses, the circumscribed parallelogram always has the very same area, namely $4ab$, where $a$ and $b$ are the semi-axes of the ellipse. Newton invokes this fact without proof, remarking only: \textit{This is demonstrated by the writers on the conic sections.}
\end{varwidth}
\end{document}
